\documentclass[12pt]{article}
\usepackage[utf8]{inputenc}
\usepackage[T1]{fontenc}
\usepackage{hyperref}
\title{Uni Language Resources for K-12 Curriculum Development}
\author{Compiled by Claude Sonnet 5 (claude-sonnet-5) for Andr\'as Kornai}
\date{}
\begin{document}
\maketitle
\section{Introduction}
Uni (also called Ramo; ISO 639-3 uni) is a Skou language of Sandaun Province, Papua New Guinea, spoken in and around Ramo village near the Indonesian border. Wikipedia gives about 670 speakers for 2000. My search found very little on Uni itself.
\section{Primary Educational Resources}
None found.
\section{Linguistic and Reference Materials}
Overview article: \url{https://en.wikipedia.org/wiki/Uni_language}. Sources differ on its classification within Sko. My search results mostly covered neighbouring Vanimo and Warapu, which I have not listed as Uni resources.
\section{Religious and Community Resources}
None found.
\section{Digital and Audio Resources}
None found.
\section{Curriculum Development Context}
The language is not known to be taught in schools, and I found no teaching materials. A curriculum would need new documentation work. Possible starting points are literature on the Sko languages and local education authorities.
\end{document}